% ============================================================================
%  Audit technique — Semantic Firewall DD
%  Cible : soumission atelier FinNLP @ EMNLP
%  Compilation : pdflatex audit_semantic_firewall.tex  (x2 pour la table)
% ============================================================================
\documentclass[11pt,a4paper]{article}

\usepackage[T1]{fontenc}
\usepackage[utf8]{inputenc}
\usepackage[french]{babel}
\usepackage{lmodern}
\usepackage[a4paper,margin=2.3cm]{geometry}
\usepackage{xcolor}
\usepackage[most]{tcolorbox}
\usepackage{listings}
\usepackage{booktabs}
\usepackage{longtable}
\usepackage{array}
\usepackage{enumitem}
\usepackage{fancyhdr}
\usepackage{amssymb}
\usepackage[hidelinks]{hyperref}

% ── Couleurs ───────────────────────────────────────────────────────────────
\definecolor{cFatal}{HTML}{B3261E}
\definecolor{cFatalBg}{HTML}{FDECEA}
\definecolor{cSerieux}{HTML}{B26A00}
\definecolor{cSerieuxBg}{HTML}{FFF4E5}
\definecolor{cMineur}{HTML}{5F6368}
\definecolor{cMineurBg}{HTML}{F1F3F4}
\definecolor{cBon}{HTML}{1E7B34}
\definecolor{cBonBg}{HTML}{EAF6EC}
\definecolor{cExp}{HTML}{1A4FA0}
\definecolor{cExpBg}{HTML}{EAF0FB}
\definecolor{cCode}{HTML}{F7F7F9}
\definecolor{cComment}{HTML}{6A737D}
\definecolor{cKw}{HTML}{A626A4}
\definecolor{cStr}{HTML}{50A14F}

% ── En-têtes ───────────────────────────────────────────────────────────────
\pagestyle{fancy}
\fancyhf{}
\fancyhead[L]{\small\itshape Audit — Semantic Firewall DD}
\fancyhead[R]{\small\itshape Cible : FinNLP @ EMNLP}
\fancyfoot[C]{\small\thepage}
\renewcommand{\headrulewidth}{0.4pt}

% ── Listings ───────────────────────────────────────────────────────────────
\lstdefinestyle{py}{
  language=Python,
  basicstyle=\ttfamily\footnotesize,
  keywordstyle=\color{cKw}\bfseries,
  commentstyle=\color{cComment}\itshape,
  stringstyle=\color{cStr},
  backgroundcolor=\color{cCode},
  frame=single,
  rulecolor=\color{cMineur!30},
  numbers=none,
  breaklines=true,
  showstringspaces=false,
  columns=fullflexible,
  xleftmargin=6pt,
  xrightmargin=6pt,
  framexleftmargin=6pt,
  aboveskip=8pt,
  belowskip=8pt,
  extendedchars=true,
  literate={é}{{\'e}}1 {è}{{\`e}}1 {ê}{{\^e}}1 {à}{{\`a}}1
           {ç}{{\c c}}1 {ô}{{\^o}}1 {î}{{\^i}}1 {û}{{\^u}}1
           {≠}{{$\neq$}}1 {→}{{$\rightarrow$}}1
}
\lstset{style=py}

% ── Boîtes de constat ──────────────────────────────────────────────────────
\newtcolorbox{bon}[1]{
  enhanced, breakable,
  colback=cBonBg, colframe=cBon, coltitle=white,
  fonttitle=\bfseries, title={#1},
  boxrule=0.8pt, arc=2pt, left=6pt, right=6pt, top=5pt, bottom=5pt
}
\newtcolorbox{fatal}[1]{
  enhanced, breakable,
  colback=cFatalBg, colframe=cFatal, coltitle=white,
  fonttitle=\bfseries, title={#1},
  boxrule=0.8pt, arc=2pt, left=6pt, right=6pt, top=5pt, bottom=5pt
}
\newtcolorbox{serieux}[1]{
  enhanced, breakable,
  colback=cSerieuxBg, colframe=cSerieux, coltitle=white,
  fonttitle=\bfseries, title={#1},
  boxrule=0.8pt, arc=2pt, left=6pt, right=6pt, top=5pt, bottom=5pt
}
\newtcolorbox{exper}[1]{
  enhanced, breakable,
  colback=cExpBg, colframe=cExp, coltitle=white,
  fonttitle=\bfseries, title={#1},
  boxrule=0.8pt, arc=2pt, left=6pt, right=6pt, top=5pt, bottom=5pt
}

\newcommand{\corr}[1]{%
  \par\vspace{4pt}%
  \noindent\textbf{\textcolor{cExp}{$\blacktriangleright$ Correction recommandée.}} #1%
}
\newcommand{\fic}[1]{\texttt{\small #1}}

% ============================================================================
\title{\vspace{-1.5cm}\textbf{Audit technique du dépôt \textit{Semantic Firewall DD}}\\[4pt]
\large Diagnostic, corrections et plan d'expériences\\[2pt]
\normalsize en vue d'une soumission à l'atelier FinNLP (EMNLP)}
\author{}
\date{30 juillet 2026}

\begin{document}
\maketitle
\thispagestyle{fancy}

\begin{tcolorbox}[colback=cMineurBg, colframe=cMineur, boxrule=0.6pt, arc=2pt]
\textbf{Comment lire ce document.} Il se lit en trois temps. La
section~\ref{sec:solide} liste ce qu'il faut \emph{garder}. Les
sections~\ref{sec:fatal} à~\ref{sec:mineur} listent ce qu'il faut
\emph{corriger}, classé par gravité, avec pour chaque point la correction à
appliquer. La section~\ref{sec:exp} donne les expériences à lancer, dans
l'ordre. La section~\ref{sec:route} résume tout sous forme de feuille de route.
\end{tcolorbox}

% ===========================================================================
\section{Résumé exécutif}

\subsection*{Le dépôt en chiffres}

\begin{center}
\begin{tabular}{@{}lr@{}}
\toprule
\textbf{Élément} & \textbf{Valeur} \\
\midrule
Lignes de Python & $\approx$ 12\,400 \\
Scripts à la racine & 36 \\
Documents de référence (corpus) & 72 \\
Fonctions de test unitaire & 1 \\
Fichier \fic{README} & absent \\
Fichier \fic{requirements.txt} & absent \\
Commits & 1 \\
\bottomrule
\end{tabular}
\end{center}

\subsection*{Le verdict en une phrase}

\begin{fatal}{Verdict}
L'idée scientifique est bonne et défendable. \textbf{L'évaluation, elle, ne
tient pas} : la métrique principale est tautologique, il n'existe aucune
vérité de terrain, et les étiquettes d'anomalie sont lues dans les noms de
fichiers. En l'état, le papier serait rejeté dès la première lecture. La bonne
nouvelle : les corrections sont identifiées, et le déblocage principal (XBRL)
coûte quelques jours de tuyauterie de données, pas une refonte.
\end{fatal}

\subsection*{L'idée à défendre}

Elle est simple et elle est solide :

\begin{center}
\begin{tcolorbox}[colback=white, colframe=cExp, boxrule=1pt, arc=2pt, width=0.93\textwidth]
\centering
\textbf{Les identités comptables sont des détecteurs d'erreur gratuits,
qui ne demandent aucune annotation.}\\[4pt]
\small
Si \textsc{ebitda} $\neq$ \textsc{ebit} $+$ \textsc{d\&a}, si Actif $\neq$
Passif, si prix par action $\neq$ valorisation post-money / nombre d'actions,
alors l'extraction du LLM est fausse --- \emph{et on le sait sans disposer de
la bonne réponse}.
\end{tcolorbox}
\end{center}

C'est exactement le type de contribution qu'un atelier FinNLP accueille. Tout
le travail consiste maintenant à le \emph{mesurer honnêtement}.

% ===========================================================================
\section{Ce qui est solide --- à garder}
\label{sec:solide}

\begin{bon}{1. La thèse centrale est publiable}
Les contrôles déterministes de \fic{dd\_base.py} (lignes 209--464) implémentent
proprement quatre identités comptables dures : équilibre du bilan,
décomposition de l'actif, cohérence \textsc{ebitda}/\textsc{ebit}/\textsc{d\&a},
et cohérence de la table de capitalisation. C'est le c\oe ur du papier.
\end{bon}

\begin{bon}{2. Le parseur de montants \fic{DDTaxonomy.\_f} est excellent}
\fic{dd\_base.py}, lignes 212--257. Il gère les négatifs entre parenthèses de la
comptabilité anglo-saxonne, l'ambiguïté décimale française/anglaise
($1.234{,}56$ contre $1{,}234.56$), les multiplicateurs textuels
(\textit{million}, \textit{bn}, \textit{mm}), les espaces insécables. C'est le
genre de détail ingrat que les relecteurs de NLP financier respectent. À
décrire dans le papier, en un paragraphe.
\end{bon}

\begin{bon}{3. La conception bilingue évite un piège réel}
\fic{semantic\_monitor.py}, lignes 283--332 : l'IDF et les centroïdes sont
calculés \emph{par langue}. Un 10-K anglais n'est donc jamais pénalisé face à un
corpus français. Beaucoup se trompent sur ce point. La liste de mots vides
propre au \textit{boilerplate} des 10-K (lignes 60--68) est également un bon
choix de domaine.
\end{bon}

\begin{bon}{4. L'interprétabilité est déjà là, gratuitement}
\fic{\_jsd\_contributors} (\fic{semantic\_monitor.py}, lignes 137--173) décompose
le score de dérive mot par mot, avec les ratios de sur- et
sous-représentation. C'est une figure du papier, et elle ne coûte rien de plus.
\end{bon}

\begin{bon}{5. Les cas difficiles ont été choisis avec discernement}
La liste des 49 documents (\fic{benchmark\_compare.py}, lignes 60--255) documente
un \emph{défi} par dépôt : banques sans ligne « Net sales », exercices fiscaux
décalés (juin, avril), résultat opérationnel de segment confondu avec le
consolidé, amortissement de contenu de Netflix employé comme approximation de
\textsc{d\&a}, épuisement (\textit{depletion}) de Chevron. Cette taxonomie
sectorielle des modes de défaillance est \textbf{une contribution en soi} --- et
elle est aujourd'hui enterrée dans des commentaires.
\end{bon}

\begin{bon}{6. Coût et latence sont déjà instrumentés}
Le champ \fic{time\_s} est enregistré par méthode, et le \textit{runner} sait
reprendre un calcul interrompu. C'est pragmatique et cela alimentera directement
le tableau coût/précision du papier.
\end{bon}

% ===========================================================================
\section{Défauts bloquants --- rédhibitoires en relecture}
\label{sec:fatal}

\begin{fatal}{D1. La métrique principale est tautologique}
C'est le point qui, seul, condamne le papier.

Dans le pare-feu M4, \fic{pipeline.py} lignes 308--313 :
\begin{lstlisting}
if ebit_v != 0 and da_v > 0:
    computed_ebitda = ebit_v + da_v
    if ebitda_v == 0 or abs(computed_ebitda - ebitda_v)/max(abs(ebitda_v),1) > 0.15:
        ebitda_v = computed_ebitda   # on ECRASE la valeur
\end{lstlisting}

Et la métrique sur laquelle M4 est noté, \fic{benchmark\_compare.py}
lignes 440--443 :
\begin{lstlisting}
err_pct = abs(ebitda - (ebit + da)) / expected * 100
coherence = err_pct <= 15.0
\end{lstlisting}

M4 \textbf{réécrit la valeur avec la formule exacte que la métrique vérifie
ensuite, à la même tolérance de 15\,\%}. Sa « cohérence \textsc{ebitda} » vaut
donc 100\,\% \emph{par construction}, pas par mérite. Le même mécanisme rend le
\fic{score\_confiance} non informatif, puisque \fic{check\_ebitda\_consistency}
tourne après la correction.

\corr{Séparer strictement le \emph{détecteur} du \emph{correcteur}. Mesurer les
violations d'identité \textbf{avant} toute auto-correction, puis évaluer ces
violations contre une vérité de terrain externe (voir E0 et E1). Le correcteur
devient un second système, évalué séparément sur le gain de précision qu'il
apporte réellement.}
\end{fatal}

\begin{fatal}{D2. Il n'existe aucune vérité de terrain}
Le \fic{recall\_pct} (\fic{benchmark\_compare.py}, ligne 430) est la
\emph{proportion de champs non nuls} --- c'est-à-dire un taux d'extraction, pas
une justesse. À aucun moment le \textit{benchmark} 49 documents ne compare un
nombre extrait à un nombre vrai.

Pire : M4 dispose de repli par expressions régulières et de dérivations
\textsc{ebit}/\textsc{ebitda} (\fic{pipeline.py}, lignes 223--318) qui
\emph{remplissent les champs vides}. Son « rappel » dépasse donc celui de M2 par
construction également.

\textbf{Les deux axes du \textit{benchmark} sont truqués en faveur de M4, et
aucun des deux ne mesure la justesse.}

\corr{Construire la vérité de terrain depuis XBRL (expérience E0). C'est
gratuit, exact et disponible à l'échelle de plusieurs centaines de dépôts.
Redéfinir ensuite le rappel comme « proportion de champs \emph{corrects} »,
avec tolérance explicite.}
\end{fatal}

\begin{fatal}{D3. Les étiquettes d'anomalie sont lues dans les noms de fichiers}
\fic{dd\_eval.py} ligne 399 et \fic{run\_all\_eval.py} ligne 110 :
\begin{lstlisting}
gt_positive = any(kw in nom for kw in ("erreur", "atypique"))
\end{lstlisting}

Précision, rappel et F1 sont donc calculés sur 15 documents dont l'auteur a
encodé l'étiquette dans le nom du fichier, et dont il a écrit le contenu pour
déclencher les contrôles. C'est un test unitaire auto-réalisateur présenté comme
une évaluation. Après exclusions, la population effective tombe à
$\approx$ 14 documents : bien trop peu pour publier un F1, quelle que soit la
méthode d'étiquetage.

\corr{Supprimer entièrement cette matrice de confusion. La remplacer par des
étiquettes dérivées de XBRL sur des dépôts réels (E0), et ne publier des
métriques de classification que sur plusieurs centaines de documents.}
\end{fatal}

\begin{fatal}{D4. La baseline « regex » est un homme de paille}
La méthode M1 utilise, dans \fic{benchmark\_compare.py} lignes 270--292 :
\begin{lstlisting}
r'Net\s+sales?\s+([\d,]+)'      # exige le nombre ACCOLE au libelle
\end{lstlisting}
tandis que le repli \emph{interne} de M4, pour le même champ
(\fic{pipeline.py}, ligne 225), utilise :
\begin{lstlisting}
r'Net\s+sales[^\n]{0,60}?([\d,]{4,})'   # absorbe les points de conduite
\end{lstlisting}

\textbf{La baseline a reçu les motifs cassés, le système les motifs qui
marchent.} Les échecs de M1 sont un artefact de l'ablation, pas une propriété
des expressions régulières. Un relecteur qui ouvre les deux fichiers le voit
immédiatement.

\corr{Donner à M1 \emph{exactement} les mêmes motifs que le repli interne de M4.
Toute autre configuration est indéfendable. Si M1 remonte fortement, c'est une
information — pas un problème.}
\end{fatal}

\begin{fatal}{D5. La provenance annoncée ne correspond pas aux données}
Deux cas, tous deux à corriger avant toute diffusion.

\textbf{(a)} Les fichiers \fic{samples/fnspid/*.txt} sont des bilans écrits à la
main, dont l'en-tête indique
\fic{Source: FNSPID dataset --- SEC 10-K filing}. Or FNSPID (Dong \textit{et
al.}, 2024) est un jeu de données de \emph{presse financière et de cours} : il ne
contient aucun bilan de 10-K.

\textbf{(b)} \fic{generate\_finnlp2020\_drift\_report.py} revendique le cadre de
Montariol \textit{et al.} (FinNLP 2020) mais évalue sur \textbf{six documents
écrits par l'auteur, avec \fic{expected\_alert} codé en dur à côté de chacun}
(lignes 28--71). Écrire des documents bourrés de vocabulaire hors corpus, puis
constater qu'une mesure de divergence lexicale se déclenche dessus, n'a aucune
validité externe.

Soumettre une provenance de jeu de données fabriquée relève de l'\textbf{intégrité
scientifique}, pas seulement de la faiblesse méthodologique.

\corr{Corriger les en-têtes de \fic{samples/fnspid/} \textbf{immédiatement} : ces
fichiers sont synthétiques et doivent le dire. Supprimer les deux études de
dérive synthétiques et les remplacer par de vrais décalages de distribution
(expérience E4).}
\end{fatal}

% ===========================================================================
\section{Défauts sérieux}
\label{sec:serieux}

\begin{serieux}{D6. Le composant le plus original est du code mort}
\fic{check\_transcription\_divergence} (\fic{dd\_eval.py}, lignes 180--230)
vérifie que chaque montant extrait apparaît près de sa ligne d'ancrage comptable
--- autrement dit, c'est un \textbf{détecteur d'hallucination par ancrage
textuel}. Il n'est jamais appelé par \fic{pipeline.py}.

Il est de plus auto-neutralisé : les lignes 215--218 se rabattent sur une
recherche dans \emph{tout} le document quand le test localisé échoue, ce qui
détruit précisément le signal que le test existe pour mesurer.

C'est le composant à plus forte valeur pour le papier, et il est à la fois
débranché et désamorcé.

\corr{Le brancher dans \fic{pipeline.py}, \textbf{supprimer le repli sur le
document entier}, puis balayer la fenêtre d'ancrage ($\pm 0$, $\pm 2$, $\pm 5$,
$\pm 10$ lignes). Voir E5.}
\end{serieux}

\begin{serieux}{D7. Le seuil JSD de 0,55 n'est ni motivé ni centralisé}
Aucun balayage, aucune courbe ROC, aucune dérivation. La valeur est de plus
recopiée en dur dans huit endroits, y compris comme chaîne de caractères dans
les générateurs de rapport : elle peut donc se désynchroniser silencieusement de
\fic{JSD\_ALERT\_THRESHOLD}.

\corr{Une seule source de vérité : importer la constante, ne jamais retaper
« 0.55 ». Puis justifier le seuil par un balayage et une AUROC (E4).}
\end{serieux}

\begin{serieux}{D8. Le corpus est trop petit, et déséquilibré}
Décompte réel par couple (type, langue) :

\begin{center}
\begin{tabular}{@{}lrr@{}}
\toprule
\textbf{Type [langue]} & \textbf{Docs} & \textbf{Termes} \\
\midrule
\fic{compte\_resultat [en]} & 16 & 2\,945 \\
\fic{compte\_resultat [fr]} & 12 & 97 \\
\fic{bilan [fr]}            & 14 & 140 \\
\fic{bilan [en]}            & 3  & 385 \\
\fic{captable [fr]}         & 9  & 79 \\
\fic{captable [en]}         & \textbf{1} & 57 \\
\bottomrule
\end{tabular}
\end{center}

La JSD est un \emph{minimum} sur les références document par document
(\fic{semantic\_monitor.py}, ligne 375) : avec $n = 1$, la distribution de
référence \emph{est} un seul document. Surtout, 2\,945 termes contre 97 signifie
que les deux « langues » ne sont pas des conditions comparables : on compare du
texte intégral de 10-K à de courts états synthétiques. \textbf{Toute comparaison
FR/EN mesure ici la longueur et le registre, pas la langue.}

\corr{Apparier longueur et registre entre conditions. Viser au minimum quelques
dizaines de documents réels par couple (type, langue), et retirer toute
conclusion FR/EN jusque-là.}
\end{serieux}

\begin{serieux}{D9. Trois scripts ne peuvent pas s'exécuter}
\fic{generate\_finnlp2020\_drift\_report.py} (l.~18--19),
\fic{generate\_fnspid\_drift\_report.py} (l.~13--14) et
\fic{generate\_samples\_sec\_report.py} (l.~12--13) contiennent en dur :
\begin{lstlisting}
os.chdir(r'c:\Users\arthus.de.chaudenay\Projet')
\end{lstlisting}
De plus, l'appel \fic{SemanticMonitor().fit('samples/dd','samples')} y omet
\fic{samples/references/*}. Ces rapports utilisent donc un corpus
\emph{différent et plus petit} que celui de l'API : \textbf{les chiffres des PDF
ne sont pas ceux que le système produit}.

\corr{Supprimer ces trois scripts, ou les réécrire avec des chemins relatifs et
le \emph{même} appel \fic{fit()} que \fic{get\_monitor()}. En l'état, ils
donnent une image fausse du système.}
\end{serieux}

\begin{serieux}{D10. Le modèle de base est trop faible pour soutenir la thèse}
Tout tourne sur \fic{meta/llama-3.1-8b-instruct}. Or « un échafaudage améliore
un modèle 8B en zéro-\textit{shot} » n'est pas un résultat en 2026. Le
relecteur demandera si un modèle de pointe avec décodage JSON contraint ne
résout pas simplement le problème.

\corr{Ajouter au moins un modèle de pointe, et un décodage contraint par
schéma JSON. Si la couche symbolique n'aide que le 8B, c'est \emph{aussi} un
résultat publiable : « petit modèle $+$ contraintes symboliques bon marché
$\approx$ modèle de pointe pour $1/20^{\text{e}}$ du coût » est un très bon
argument d'atelier, et le tableau de coûts est déjà instrumenté.}
\end{serieux}

\begin{serieux}{D11. Le dépôt n'est pas reproductible}
Pas de \fic{README}, pas de \fic{requirements.txt}, aucune version épinglée,
aucune graine aléatoire, aucun instantané de données figé.
\fic{samples/real\_world/} est dans le \fic{.gitignore} --- \textbf{le corpus de
49 documents est donc absent du dépôt} --- et \fic{output/} aussi, donc aucun
résultat n'est versionné. \textbf{Personne ne peut aujourd'hui reproduire un
seul chiffre.}

\corr{Voir la liste d'actions de la section~\ref{sec:ingenierie}.}
\end{serieux}

\begin{serieux}{D12. Le dépôt est un produit, pas un artefact de recherche}
\fic{api.py} pèse 72\,ko, dont environ 830 lignes de HTML en ligne, auxquelles
s'ajoutent les connecteurs Notion, Airtable, Teams et SharePoint. Rien de tout
cela n'a sa place dans une soumission, et tout cela dilue la relecture.

\corr{Extraire \fic{api.py}, \fic{connectors/} et les générateurs de rapports
HTML/PDF dans un dépôt ou une branche \fic{product/} séparée. L'artefact soumis
ne contient que : extraction, validation, évaluation.}
\end{serieux}

% ===========================================================================
\section{Défauts mineurs}
\label{sec:mineur}

\begin{itemize}[leftmargin=*, itemsep=4pt]
\item \textbf{Détection de langue.} \fic{detect\_language} renvoie \fic{en} en
cas d'égalité, y compris sur une chaîne vide (\fic{semantic\_monitor.py},
l.~99). La classe de caractères inclut les accents mais oublie \fic{\oe} et
\fic{ÿ}.

\item \textbf{Règles de correction non idempotentes.} Le cas~3
(\fic{pipeline.py}, l.~289--294) modifie \fic{ebit} à partir de
$\textsc{ebitda} - \textsc{d\&a}$, puis la règle générale peut recalculer
\fic{ebitda} depuis ce même \fic{ebit}. L'ordre compte, et rien ne le teste.

\item \textbf{Les pertes sont silencieusement ignorées.} Les replis exigent
\fic{\_v >= 100} et les corrections testent \fic{> 0} : les dépôts déficitaires
sont écartés sans trace --- alors que la perte opérationnelle de Boeing figure
explicitement dans la liste des 49 documents.

\item \textbf{Troncature invisible.} \fic{\_certify\_dd} coupe à 200\,k
caractères (\fic{pipeline.py}, l.~176) après que \fic{\_prepare\_text} a déjà
plafonné à 400\,k. La troncature n'est jamais journalisée : la perte de données
est donc invisible dans les résultats.

\item \textbf{Latence faussée.} \fic{time.sleep(1.0)} est appelé à l'intérieur
de \fic{call\_dd\_extraction} (\fic{dd\_eval.py}, l.~61) et avant chaque méthode
mesurée (\fic{benchmark\_compare.py}, l.~526) : cela gonfle la latence des
méthodes que l'on prétend chronométrer.

\item \textbf{Identifiants mélangés FR/EN} dans tout le code : les relecteurs
qui lisent le code trébucheront.

\item \textbf{Provenance à vérifier.} Le fichier\\
\fic{samples/references/bilan/bilan\_devoteam\_2022.txt}\\
contient les comptes
consolidés de votre employeur. Les chiffres 2022 ont été publiés, donc c'est
très probablement sans risque --- mais vérifiez la concordance avec le rapport
annuel public avant diffusion d'un artefact ouvert.
\end{itemize}

% ===========================================================================
\section{Plan d'expériences}
\label{sec:exp}

Classées par gain de probabilité d'acceptation par unité de travail.

\subsection{Palier 0 --- sans ceci, il n'y a pas de papier}

\begin{exper}{E0. Construire la vérité de terrain depuis XBRL --- \textit{le} déblocage}
Les jeux \textit{Financial Statement Data Sets} et l'API \fic{companyfacts} de
la SEC fournissent, pour chaque dépôt, les valeurs balisées exactes et lisibles
par machine :
\fic{us-gaap:Revenues} (ou
\fic{RevenueFromContractWithCustomerExcludingAssessedTax}),
\fic{OperatingIncomeLoss},
\fic{DepreciationDepletionAndAmortization}, \fic{NetIncomeLoss},
\fic{Assets}, \fic{LiabilitiesAndStockholdersEquity},
\fic{StockholdersEquity}.

Cela donne une \textbf{vérité de terrain exacte sur 500 à 2\,000 dépôts, pour un
coût d'annotation nul}. Appariement sur CIK $+$ période fiscale $+$ échelle
d'unité.

\textbf{À faire en premier : tout le reste en dépend.} Chaque métrique fausse du
dépôt devient alors une vraie métrique — justesse en correspondance exacte ou
tolérée, par champ, par secteur, par méthode.
\end{exper}

\begin{exper}{E1. Séparer détecteur et correcteur, puis évaluer le détecteur}
Geler l'extraction. Calculer les violations d'identité \textbf{avant} toute
auto-correction. Puis, avec les étiquettes E0, poser la véritable question de
recherche :

\begin{center}
\itshape Une violation d'identité comptable prédit-elle qu'une valeur est fausse ?
\end{center}

Rapporter ROC-AUC, PR-AUC, précision@k, et la calibration de
\fic{score\_confiance} contre les étiquettes d'erreur dérivées de XBRL.

\textbf{C'est cela, votre papier.} Un détecteur d'erreur non supervisé, sans
étiquettes, avec une AUROC mesurée sur de vrais dépôts, est une contribution
FinNLP propre. La boucle actuelle « je corrige, puis je vérifie avec la même
formule » n'en est pas une.
\end{exper}

\begin{exper}{E2. Réparer les baselines, puis publier l'échelle complète}
Donner d'abord à M1 les \emph{mêmes} motifs que le repli interne de M4. Puis :

\begin{center}
\begin{tabular}{@{}ll@{}}
\toprule
\textbf{Réf.} & \textbf{Méthode} \\
\midrule
B0 & Oracle XBRL (borne supérieure) \\
B1 & Expressions régulières, réglées de bonne foi \\
B2 & LLM zéro-\textit{shot}, décodage contraint par schéma JSON \\
B3 & B2 $+$ auto-cohérence sur $k=5$, vote majoritaire \\
B4 & B2 $+$ auto-vérification CoT (votre M3 actuel) \\
B5 & B2 $+$ couche d'identités déterministes (M4, \emph{correcteur désactivé}) \\
B6 & B5 $+$ correction, gains attribués règle par règle \\
\bottomrule
\end{tabular}
\end{center}

Deux niveaux de modèle au minimum : un 8B ouvert et un modèle de pointe.
\end{exper}

\begin{exper}{E3. B4 contre B5 : l'auto-vérification du LLM corrige-t-elle l'arithmétique ?}
Votre invite \fic{\_COT\_VERIFY} (\fic{benchmark\_compare.py}, l.~356--366)
demande au modèle de revérifier une identité \emph{sur sa propre sortie} : c'est
un biais de confirmation de manuel. Mesurez-le contre les étiquettes E0.

\textbf{Un résultat négatif propre ici est très publiable en atelier}, et il
motive directement la couche déterministe. Publier la matrice de confusion
« le CoT a changé la valeur » $\times$ « la valeur est devenue correcte ».
\end{exper}

\subsection{Palier 1 --- fait passer d'un papier mince à un papier solide}

\begin{exper}{E4. Remplacer l'étude de dérive synthétique par un vrai décalage de distribution}
Supprimer entièrement les évaluations \fic{finnlp2020/} et \fic{fnspid/}
écrites à la main. Substituer :

\begin{itemize}[leftmargin=*, itemsep=2pt]
\item \textbf{Diachronique} : texte réel des \textit{Item 7/8} de 10-K, corpus
2015--2019 contre dépôts 2020--2021. Le glissement lexical lié au COVID est
réellement présent, dans de vrais documents.
\item \textbf{Synchronique} : retirer un secteur SIC entier du corpus de
référence, tester dessus. L'écart de vocabulaire banque/technologie est réel et
n'a pas besoin d'être fabriqué.
\item \textbf{Hors distribution franc} : texte non financier de longueur
appariée.
\end{itemize}

Puis : \textbf{balayer le seuil JSD et publier une AUROC} au lieu d'affirmer
0,55. Baselines : cosinus TF-IDF au centroïde, perplexité d'un petit modèle de
langue, MMD sur plongements de phrases. Si la JSD perd contre la perplexité, le
dire ; si elle égalise en AUROC et gagne en interprétabilité, c'est
\emph{cela} l'argument, correctement établi.
\end{exper}

\begin{exper}{E5. Brancher et évaluer le contrôle d'ancrage --- votre meilleur actif inutilisé}
Activer \fic{check\_transcription\_divergence} dans \fic{pipeline.py},
\textbf{supprimer le repli sur le document entier} (\fic{dd\_eval.py},
l.~215--218), et balayer la fenêtre d'ancrage.

Évaluer comme détecteur d'hallucination contre E0 : quand la valeur extraite
diffère de XBRL, le contrôle d'ancrage se déclenche-t-il ? Vous obtenez ainsi un
\textbf{second signal d'erreur sans étiquettes, indépendant du premier}, et la
combinaison « identité arithmétique $\vee$ ancrage textuel » constitue un
système plus intéressant que chacun pris isolément. Publier les deux séparément,
puis combinés.
\end{exper}

\begin{exper}{E6. Ablation des règles de correction}
Huit règles réglées à la main dans \fic{pipeline.py}, l.~272--325.
Retrait d'une règle à la fois, mesuré contre E0. Il est probable que deux ou
trois portent tout le gain, et qu'une ou deux \emph{dégradent} le résultat --- le
cas~3, qui modifie \fic{ebit}, est le premier suspect. Tester aussi l'idempotence
et la sensibilité à l'ordre.

Annoncer soi-même « 3 de mes 8 heuristiques étaient nuisibles » vaut beaucoup de
crédibilité en relecture.
\end{exper}

\begin{exper}{E7. Analyse d'erreur stratifiée par secteur}
La taxonomie sectorielle et les hypothèses de défaillance par document existent
déjà (\fic{benchmark\_compare.py}, l.~60--255). Avec les étiquettes E0,
promouvez-les du commentaire au tableau principal du papier : justesse
d'extraction $\times$ secteur $\times$ champ, avec confirmation ou réfutation du
mode de défaillance prédit, cellule par cellule.

\textbf{Cette taxonomie d'erreurs est la contribution autonome la plus
défendable du dépôt}, et c'est la partie que les relecteurs citeront.
\end{exper}

\subsection{Palier 2 --- si le temps le permet}

\begin{itemize}[leftmargin=*, itemsep=4pt]
\item \textbf{E8. Robustesse.} Injection de bruit d'OCR et d'extraction PDF,
corruption au niveau caractère, perturbations d'aplatissement de tableaux
$\rightarrow$ courbes de dégradation par méthode. Parle directement aux
praticiens.
\item \textbf{E9. Translingue.} De vrais dépôts français (BALO, rapports annuels
IFRS Euronext), et non des documents synthétiques. La thèse bilingue aura alors
un support.
\item \textbf{E10. Front de Pareto coût / latence / justesse.} \fic{time\_s} est
déjà journalisé ; retirer d'abord la contamination par \fic{time.sleep}.
\end{itemize}

% ===========================================================================
\section{Chantier d'ingénierie --- avant de lancer quoi que ce soit}
\label{sec:ingenierie}

\begin{enumerate}[leftmargin=*, itemsep=4pt]
\item Extraire \fic{api.py}, \fic{connectors/} et tous les
\fic{generate\_*\_report.py} dans un dépôt ou une branche \fic{product/}
séparée. L'artefact de soumission ne livre qu'extraction, validation et
évaluation.
\item Ajouter \fic{README.md}, \fic{requirements.txt} avec versions épinglées,
un \fic{Makefile} qui reproduit chaque tableau, et des graines aléatoires fixes.
\item Versionner un instantané de données immuable : CIK $+$ numéros d'accession
$+$ identifiants de faits XBRL. Ne pas committer les PDF --- committer un script
de récupération et un manifeste avec sommes de contrôle.
\item Supprimer ou réparer les trois scripts en \fic{os.chdir(r'c:\textbackslash
Users\textbackslash arthus...')} : ils donnent aujourd'hui une image fausse du
corpus réellement utilisé.
\item Corriger les en-têtes de provenance de \fic{samples/fnspid/}
\textbf{maintenant}, avant toute diffusion.
\item Une seule source de vérité pour \fic{JSD\_ALERT\_THRESHOLD} : l'importer,
ne jamais retaper « 0.55 ».
\item Écrire de vrais tests sur \fic{DDTaxonomy.\_f} (votre meilleur code, avec
une couverture quasi nulle) et sur l'ordre des règles de correction.
\end{enumerate}

% ===========================================================================
\section{Feuille de route et verdict}
\label{sec:route}

\subsection*{Synthèse défaut $\rightarrow$ correction}

{\small
\begin{longtable}{@{}p{0.055\textwidth}p{0.36\textwidth}p{0.44\textwidth}p{0.07\textwidth}@{}}
\toprule
\textbf{Réf.} & \textbf{Défaut} & \textbf{Correction} & \textbf{Exp.} \\
\midrule
\endfirsthead
\toprule
\textbf{Réf.} & \textbf{Défaut} & \textbf{Correction} & \textbf{Exp.} \\
\midrule
\endhead
D1 & Métrique tautologique & Séparer détecteur et correcteur & E0, E1 \\
D2 & Aucune vérité de terrain & Étiquettes XBRL & E0 \\
D3 & Étiquettes dans les noms de fichiers & Étiquettes XBRL sur dépôts réels & E0 \\
D4 & Baseline regex homme de paille & Motifs identiques à ceux de M4 & E2 \\
D5 & Provenance fabriquée & Corriger les en-têtes ; vrai décalage & E4 \\
D6 & Détecteur d'ancrage en code mort & Brancher, retirer le repli global & E5 \\
D7 & Seuil 0,55 arbitraire & Balayage $+$ AUROC ; constante unique & E4 \\
D8 & Corpus trop petit, déséquilibré & Apparier longueur et registre & E0, E9 \\
D9 & Trois scripts inexécutables & Chemins relatifs ou suppression & --- \\
D10 & Modèle de base trop faible & Ajouter un modèle de pointe & E2 \\
D11 & Non reproductible & README, versions, instantané, graines & --- \\
D12 & Dépôt-produit & Extraire l'API et les connecteurs & --- \\
\bottomrule
\end{longtable}
}

\subsection*{Ordre d'exécution}

\begin{center}
\begin{tcolorbox}[colback=cExpBg, colframe=cExp, boxrule=0.8pt, arc=2pt, width=0.95\textwidth]
\textbf{1.} Nettoyage d'intégrité : en-têtes \fic{fnspid}, scripts inexécutables,
constante de seuil.\\
\textbf{2.} \textbf{E0 --- XBRL.} Rien d'autre ne compte avant.\\
\textbf{3.} E1 (détecteur seul) $\rightarrow$ E2 (baselines) $\rightarrow$ E3
(CoT, résultat négatif).\\
\textbf{4.} Rédaction possible dès ici : \emph{papier court}.\\
\textbf{5.} E4, E5, E6, E7 $\rightarrow$ \emph{papier long} confortable.
\end{tcolorbox}
\end{center}

\subsection*{Verdict}

\begin{fatal}{En l'état}
\textbf{Probabilité d'acceptation : très faible.} Non parce que l'idée est
faible --- elle ne l'est pas --- mais parce que chaque chiffre mis en avant est
soit tautologique (M4 noté sur la formule qu'il applique), soit sans étiquettes
par accident (rappel $=$ taux de non-nuls), soit mesuré sur des documents que
l'auteur a écrits pour produire le résultat souhaité. Un seul relecteur lisant
\fic{pipeline.py} ligne 308 en regard de \fic{benchmark\_compare.py} ligne 440
suffit à conclure.
\end{fatal}

\begin{bon}{Avec E0 à E3 menées honnêtement}
\textbf{Réellement compétitif} pour FinNLP @ EMNLP en papier court. Vérité de
terrain XBRL à l'échelle, détecteur d'erreur sans étiquettes avec AUROC réelle,
résultat négatif propre sur l'auto-vérification des LLM, et taxonomie d'erreurs
stratifiée par secteur sur de vrais 10-K : cela forme un papier d'atelier
cohérent, honnête et utile.

Ajouter E4 à E7 en fait un papier long confortable.
\end{bon}

\vspace{6pt}
\noindent
L'action de plus fort levier est \textbf{E0}. XBRL transforme toute l'évaluation
factice en évaluation réelle pour le prix de quelques jours de tuyauterie de
données --- et rien d'autre dans cette liste ne compte avant que ce soit fait.

% ===========================================================================
\section{Addendum --- état après corrections}
\label{sec:addendum}

Les défauts D1 à D12 ont été corrigés (paquet \fic{semantic\_firewall/}, 49 tests).
Un second passage d'audit a révélé trois défauts \emph{introduits ou rendus visibles
par ces corrections elles-mêmes}. Ils sont désormais corrigés aussi, et E0--E3 puis
E5--E7 ont été exécutés sur données réelles.

\begin{fatal}{D13 --- Incompatibilité d'échelle (bloquant, corrigé)}
XBRL publie en dollars absolus ($416\,161\,000\,000$) ; les tableaux impriment des
millions ($416\,161$) ; et le prompt ordonne « ne PAS multiplier ». Conséquence :
\textbf{100\,\% des comparaisons échouaient}. Le code tournait et renvoyait des
nombres --- tous faux dans le même sens, ce qui est pire qu'un plantage : $0\,\%$ de
justesse pour \emph{toutes} les baselines se lit comme un résultat de modèle, pas
comme une erreur d'unité. L'ancienne suite de tests ne pouvait pas le voir : ses
fixtures construisaient l'extraction \emph{depuis} la vérité de terrain, donc les
unités coïncidaient par construction.

\corr{Appliquée. Nouveau module \fic{extraction/scale.py} : le multiplicateur est
inféré \textbf{depuis le texte du dépôt}, jamais depuis les étiquettes --- s'aligner
sur la puissance de 1000 la plus proche de la GT ferait fuiter la réponse dans la
prédiction et recréerait D1 ailleurs. Les dépôts sans échelle déclarée sont exclus
(13 documents), pas devinés. \fic{classify\_error()} sépare désormais
\emph{correct} / \emph{échelle} / \emph{signe} / \emph{valeur fausse} : une erreur
d'unité laisse toutes les identités exactement satisfaites, elle est donc invisible à
cette classe de méthode et exclue des étiquettes du détecteur.}
\end{fatal}

\begin{serieux}{D14 --- Décalage d'un an sur 32/46 enregistrements (corrigé)}
\fic{nearest\_fiscal\_year()} retombait sur l'exercice disponible le plus proche : tous
les documents étiquetés 2026 se résolvaient en FY2025. Un an d'écart sur le chiffre
d'affaires dépasse toute tolérance et est indiscernable d'une erreur d'extraction.

\corr{Corrigé à la racine : l'exercice est lu dans le \fic{reportDate} du dépôt
lui-même, donc document et étiquette portent sur la même période par construction.
Filet conservé : \fic{load()} exclut par défaut tout enregistrement discordant en
rapportant le décompte.}
\end{serieux}

\begin{serieux}{D15 --- Couverture par champ non publiée (corrigé)}
L'identité EBITDA n'était évaluable que sur 26 des 46 documents, sans que ce soit dit.

\corr{\fic{coverage\_report()} publie le $n$ effectif par champ et compte les champs
dérivés. Sur le nouveau corpus : de 63,5\,\% (EBITDA, dérivé) à 100\,\% (actif,
passif, capitaux propres, résultat net).}
\end{serieux}

\subsection*{Ce que ces corrections ont permis de mesurer}

Le corpus est passé à \textbf{104 dépôts 10-K réels, 55 sociétés, 11 secteurs}, avec
vérité de terrain XBRL. Deux résultats n'étaient pas atteignables avant.

\begin{bon}{Le résultat central : la couverture, pas la force de l'identité}
Même code, même corpus, même modèle :
\begin{itemize}[leftmargin=*, itemsep=1pt]
\item \textbf{Compte de résultat} --- l'EBITDA n'étant pas une ligne GAAP, le modèle
renvoie \fic{null} dans \textbf{0/69} cas, l'identité n'est \emph{jamais} calculable,
et l'AUROC de $0{,}500$ est une dégénérescence par ex-æquo, pas une mesure. Le
détecteur n'est pas imprécis, il est \emph{inapplicable}.
\item \textbf{Bilan} --- les termes sont imprimés partout, l'identité est calculable
sur \textbf{68/69}, et l'on obtient \textbf{AUROC $0{,}753$, PR-AUC $0{,}791$, et une
précision de $1{,}00$ sur les documents les mieux classés} --- directement exploitable
pour du triage, sans aucune étiquette.
\end{itemize}
Règle pratique qui en découle : bâtir la vérification sur les grandeurs que le dépôt
est \emph{obligé} de publier.
\end{bon}

\begin{bon}{Le leave-one-out mentait (E6)}
Retirer une règle à la fois donne $\Delta = 0{,}0$ pour les huit règles --- conclusion
naturelle : ce sont des décorations. L'ablation \textbf{par bloc} montre l'inverse :
retirer le bloc arithmétique coûte $-11{,}2$ pt (8B) et $-13{,}8$ pt (49B). Les deux
mesures sont justes : le jeu de règles est \emph{redondant} sous itération jusqu'au
point fixe, et le leave-one-out y est structurellement aveugle. Mon hypothèse initiale
(« 2 ou 3 règles portent tout, 1 ou 2 nuisent ») était fausse dans les deux sens.
\end{bon}

\noindent
Autres résultats : le correcteur déterministe atteint $84{,}5\,\%$ (8B) et
$91{,}2\,\%$ (49B) contre $73{,}3$ / $77{,}4$ en zéro-\textit{shot}, \textbf{pour zéro
appel supplémentaire}, battant l'auto-vérification CoT ($79{,}7$ / $89{,}9$) qui coûte
le double. L'auto-cohérence à $k{=}5$ n'apporte rien pour cinq fois le coût. Le CoT,
contrairement à mon attente, \emph{aide} (19 champs corrigés, 0 dégradé) --- mais
seulement parce que le prompt lui fournit le critère externe.

Le papier correspondant est \fic{paper/main.tex} (14 pages, format ACL). Aucun
résultat de dérive lexicale n'y est publié : le seuil de $0{,}55$ n'a toujours pas été
justifié par balayage, et le papier le déclare en Limitations plutôt que de reprendre
les anciens chiffres synthétiques.

\end{document}
